% !TEX program = xelatex
\documentclass[a4paper,12pt]{article}

% ==================== 中文支持 ====================
\usepackage[UTF8]{ctex}

% ==================== 页面设置 ====================
\usepackage{geometry}
\geometry{left=2.5cm,right=2.5cm,top=2.5cm,bottom=2.5cm}

% ==================== 数学公式 ====================
\usepackage{amsmath}
\usepackage{amsfonts}
\usepackage{amssymb}
\usepackage{bm}

% ==================== 图表 ====================
\usepackage{graphicx}
\usepackage{float}
\usepackage{caption}
\captionsetup{labelsep=quad}

% ==================== 表格 ====================
\usepackage{array}
\usepackage{multirow}
\usepackage{booktabs}
\usepackage{longtable}

% ==================== 代码 ====================
\usepackage{listings}
\usepackage{xcolor}
\lstset{
	basicstyle=\small\ttfamily,
	breaklines=true,
	frame=single,
	numbers=left,
	numberstyle=\tiny,
	keywordstyle=\color{blue},
	commentstyle=\color{green!60!black},
	stringstyle=\color{red},
	showstringspaces=false
}

% ==================== 其他 ====================
\usepackage{hyperref}
\usepackage{indentfirst}
\setlength{\parindent}{2em}

% ==================== 标题格式 ====================
\usepackage{titlesec}
\titleformat{\section}{\centering\Large\bfseries}{\zhnum{section}、}{0em}{}
\titleformat{\subsection}{\large\bfseries}{\thesubsection}{1em}{}

% ==================== 正文开始 ====================
\begin{document}
	
	% ==================== 标题页 ====================
	\begin{center}
		{\LARGE\bfseries 某商超蔬菜类商品动态定价与补货决策研究} \\[0.5cm]
		{\large 摘要} \\[0.3cm]
	\end{center}
	
	\noindent 随着生鲜市场规模的持续扩大，蔬菜零售行业的竞争也愈加激烈。本文针对蔬菜类商品的动态定价与补货决策问题建立了数学模型进行分析。\\[0.3cm]
	
	\noindent \textbf{针对问题一}，首先计算描述统计量并进行数据可视化，分析各蔬菜品类和单品蔬菜的销售量分布规律；随后引入Spearman相关系数和K-means++聚类算法进行相关性分析，将单品划分为热销、畅销、平销和滞销四大类。\\[0.3cm]
	
	\noindent \textbf{针对问题二}，第一问通过Pearson相关系数检验销售总量与成本加成定价的负线性关系，建立线性回归模型；第二问构建时间序列预测模型预测未来一周销量，并以利润最大化为目标建立优化模型，引入模拟退火算法求解，得到未来一周最大收益为5105.60元及最优补货定价策略。\\[0.3cm]
	
	\noindent \textbf{针对问题三}，以最大化收益为目标，以单品种类数量、陈列量、补货量和定价为约束建立优化模型，筛选出29种可售单品，估计2023年7月1日最大收益为1282.26元。\\[0.3cm]
	
	\noindent \textbf{针对问题四}，选取季度蔬菜丰富度、季度效应和节日因素三个指标，运用灰色关联分析模型，求得各指标与六大蔬菜品类季度平均销售量的关联度均在0.5以上，表明具有强相关性。\\[0.3cm]
	
	\noindent 最后对模型进行了评价与推广。\\[0.3cm]
	
	\noindent \textbf{关键词：} K-means++聚类；优化模型；模拟退火；灰色关联分析
	
	\newpage
	
	% ==================== 目录 ====================
	% \tableofcontents
	% \newpage
	
	% ==================== 一、问题重述 ====================
	\section{问题重述}
	
	\subsection{问题背景}
	
	蔬菜作为人民生活不可或缺的食物之一，其本身具有高频性、精致化等消费特点，市场前景广阔。但由于蔬菜具有周期性、季节性、易腐和易伤等自然属性，加之蔬菜品种众多且产地不同所导致的购入成本各异，因此合理的定价和补货决策对于商家实现成本较低、利润最大化的经营目标至关重要。
	
	\subsection{问题要求}
	
	根据题目所给附件数据信息，需要解决以下问题：
	
	\textbf{问题一：} 分析蔬菜各品类、单品销售量的分布规律和相互关系，判断不同品类或不同单品之间是否存在一定的关联关系。
	
	\textbf{问题二：} 计算各蔬菜品类的成本加成定价，分析其与销售总量之间的关系；以利润最大化为目标，制定未来一周的日补货总量和定价策略。
	
	\textbf{问题三：} 根据2023年6月24-30日的可售品种，按照单品订购量不少于2.5千克等要求，筛选出27-33个可售单品，制定7月1日的单品补货量和定价策略。
	
	\textbf{问题四：} 从季节、节日、销售形式等角度分析数据特征，整理对制定补货和定价决策有较大影响的相关数据，并给出合理意见。
	
	\newpage
	
	% ==================== 二、问题分析 ====================
	\section{问题分析}
	
	\subsection{问题一的分析}
	
	针对问题一，第一小问，首先对附件2中的无效数据进行剔除处理，然后引入均值、最大值、最小值、中位数、标准差、偏度系数、峰度系数等描述统计量，并通过数据可视化分析，直观观测各品类蔬菜和单品蔬菜的销售情况及趋势，由此得出其分布规律。
	
	第二小问，引入Spearman相关系数进行各蔬菜品类销售量的相关性分析；采用K-means++聚类算法，以单品的总销售量、每日最大销售量和日均销售量为指标，将单品划分为热销、畅销、平销和滞销四大类。
	
	\subsection{问题二的分析}
	
	针对问题二，第一问先根据成本加成定价公式计算各蔬菜品类每日定价，引入Pearson相关系数判定销售总量与成本加成定价的线性关系，建立线性回归模型。
	
	第二问针对不同品类蔬菜建立时间序列预测模型，以前半年日销售量为原始数据预测未来一周日销售量，建立以利润最大化为目标的优化模型，引入模拟退火算法求解最优补货和定价策略。
	
	\subsection{问题三的分析}
	
	针对问题三，本题为优化问题中的线性规划问题，基于问题二构建的模拟退火模型。首先根据题意筛选可用数据范围，确定目标函数及约束条件，求出最优单品补货量和定价策略。
	
	\subsection{问题四的分析}
	
	针对问题四，筛选出季度各蔬菜品类的单品蔬菜销售种数、节气期间特征单品蔬菜销量和节日期间特征单品蔬菜销量，分别代表季度蔬菜丰富度、季度效应和节日三个因素，建立灰色关联分析模型，求出各指标对销售量的关联度。
	
	\newpage
	
	% ==================== 三、模型假设 ====================
	\section{模型假设}
	
	\begin{enumerate}
		\item 假设蔬菜的历史销售数据能够代表未来的销售趋势，本文使用历史销售数据来预测未来的销售。
		\item 假设蔬菜的损耗是固定的，本文使用损耗率来计算蔬菜的单位成本。
		\item 假设商超的销售空间在一段时间内是恒定的。
	\end{enumerate}
	
	\newpage
	
	% ==================== 四、符号说明 ====================
	\section{符号说明}
	
	\begin{table}[H]
		\centering
		\begin{tabular}{c|c|c}
			\hline
			\textbf{符号} & \textbf{含义} & \textbf{单位} \\
			\hline
			$x_{ik}$ & 第$i$种单品第$k$个指标 & kg \\
			$c_{jk}$ & 第$j$个聚类中心第$k$个指标 & -- \\
			$P_i$ & 第$i$种蔬菜品类的定价 & 元/kg \\
			$S_{ij}$ & 第$i$品类第$j$天的销量 & kg \\
			$R_{ij}$ & 第$i$品类第$j$天的补货量 & kg \\
			$B_i$ & 第$i$品类的单位成本 & 元/kg \\
			$L_i$ & 第$i$品类的损耗率 & $\%$ \\
			$y_i$ & 0-1变量，是否采购第$i$种单品 & -- \\
			\hline
		\end{tabular}
	\end{table}
	
	\newpage
	
	% ==================== 五、数据侧写 ====================
	\section{数据侧写}
	
	为更好地理解和分析数据，本文首先对数据进行多角度探索性分析。
	
	\subsection{六大品类蔬菜的品种丰富度}
	
	花叶类蔬菜占品种数量的40\%，品种丰富；花菜类蔬菜的品种数量只占2\%，选择范围较窄。
	
	\subsection{六大品类蔬菜季度平均销售量}
	
	花叶类较其他五类蔬菜全年销售量较高。花叶类、辣椒类、食用菌和水生根茎类销售量峰值的出现规律都呈现出周期性，大都在每年的一、三、四季度出现峰值。
	
	\subsection{六大品类蔬菜损耗率}
	
	花菜类、水生根茎类以及花叶类的蔬菜损耗率都在10\%以上，损耗率较高，较难保存。茄类的损耗率大约为7\%，较容易保存。
	
	\newpage
	
	% ==================== 六、问题一模型的建立与求解 ====================
	\section{问题一模型的建立与求解}
	
	\subsection{数据预处理}
	
	在分析蔬菜各品类及单品的销售量分布情况时，部分商品存在退货的情况，销售单价显示为负值，不能反映实际的销售量情况，属于无效数据，故对退货的数据进行剔除处理。
	
	\subsection{第一问：通过描述统计分析分布规律}
	
	偏度系数计算公式如下：
	
	\begin{equation}
		S_k = E\left[\left(\frac{x-\mu}{\sigma}\right)^3\right] = \frac{\mu^3}{\sigma^3}
		\label{eq:skewness}
	\end{equation}
	
	峰度系数计算公式如下：
	
	\begin{equation}
		K_u = E\left[\left(\frac{X-\mu}{\sigma}\right)^4\right] = \frac{\mu^4}{\sigma^4}
		\label{eq:kurtosis}
	\end{equation}
	
	\begin{table}[H]
		\centering
		\caption{蔬菜六大品类日销售量统计学指标表}
		\begin{tabular}{c|c|c|c|c|c|c}
			\hline
			\textbf{指标} & \textbf{花菜类} & \textbf{食用菌} & \textbf{花叶类} & \textbf{辣椒类} & \textbf{茄类} & \textbf{水生根茎类} \\
			\hline
			最小值 & 0.00 & 0.00 & 0.00 & 0.00 & 0.01 & 0.00 \\
			最大值 & 186.16 & 511.14 & 1265.47 & 604.23 & 118.93 & 296.79 \\
			均值 & 38.16 & 69.18 & 181.42 & 83.69 & 20.50 & 37.08 \\
			偏度 & 1.47 & 2.95 & 2.71 & 3.06 & 1.55 & 2.48 \\
			峰度 & 7.03 & 19.85 & 27.36 & 20.78 & 8.15 & 14.95 \\
			标准差 & 22.90 & 48.20 & 87.65 & 53.82 & 13.57 & 31.43 \\
			\hline
		\end{tabular}
	\end{table}
	
	由表可知，花叶类和辣椒类的均值较高，标准差也相对较大，说明其销售量的波动性较大；茄类的销售量均值较低，标准差也相对较小，说明其销售量稳定。
	
	\subsection{第二问：建立模型判断关联关系}
	
	Spearman相关系数计算公式如下：
	
	\begin{equation}
		r_s = 1 - \frac{6\sum_{i=1}^{N}d_i^2}{N(N^2-1)}
		\label{eq:spearman}
	\end{equation}
	
	\begin{table}[H]
		\centering
		\caption{蔬菜六大品类的Spearman相关系数表}
		\begin{tabular}{c|c|c|c|c|c|c}
			\hline
			& \textbf{花菜类} & \textbf{食用菌} & \textbf{花叶类} & \textbf{辣椒类} & \textbf{茄类} & \textbf{水生根茎类} \\
			\hline
			花菜类 & 1.00 & 0.48 & 0.64 & 0.45 & 0.21 & 0.41 \\
			食用菌 & 0.48 & 1.00 & 0.61 & 0.55 & -0.08 & 0.62 \\
			花叶类 & 0.64 & 0.61 & 1.00 & 0.74 & 0.16 & 0.47 \\
			辣椒类 & 0.45 & 0.55 & 0.74 & 1.00 & 0.30 & 0.57 \\
			茄类 & 0.21 & -0.08 & 0.16 & 0.30 & 1.00 & 0.34 \\
			水生根茎类 & 0.41 & 0.62 & 0.47 & 0.57 & 0.34 & 1.00 \\
			\hline
		\end{tabular}
	\end{table}
	
	K-means++聚类结果如表所示：
	
	\begin{table}[H]
		\centering
		\caption{聚类中心}
		\begin{tabular}{c|c|c|c}
			\hline
			\textbf{类别} & \textbf{总销售量} & \textbf{每日最大销售量} & \textbf{日均销售量} \\
			\hline
			滞销 & 438.85 & 11.94 & 0.40 \\
			畅销 & 14215.59 & 42.04 & 12.98 \\
			平销 & 6087.95 & 22.00 & 5.56 \\
			热销 & 28444.36 & 200.54 & 25.98 \\
			\hline
		\end{tabular}
	\end{table}
	
	\newpage
	
	% ==================== 七、问题二模型的建立与求解 ====================
	\section{问题二模型的建立与求解}
	
	\subsection{第一问：销售总量与成本加成定价的关系}
	
	Pearson相关系数计算公式如下：
	
	\begin{equation}
		r = \frac{\sum_{i=1}^{n}(X_i-\bar{X})(Y_i-\bar{Y})}{\sqrt{\sum_{i=1}^{n}(X_i-\bar{X})^2}\sqrt{\sum_{i=1}^{n}(Y_i-\bar{Y})^2}}
		\label{eq:pearson}
	\end{equation}
	
	\begin{table}[H]
		\centering
		\caption{Pearson相关系数}
		\begin{tabular}{c|c|c|c|c|c|c}
			\hline
			\textbf{品类} & 花菜类 & 食用菌 & 花叶类 & 辣椒类 & 茄类 & 水生根茎类 \\
			\hline
			\textbf{相关系数} & -0.30 & -0.28 & -0.27 & -0.26 & -0.17 & -0.21 \\
			\hline
		\end{tabular}
	\end{table}
	
	线性回归模型：
	
	\begin{equation}
		cpp = \beta_0 + \beta_1 cs
		\label{eq:regression}
	\end{equation}
	
	\subsection{第二问：优化模型制定补货和定价策略}
	
	以利润最大化为目标函数：
	
	\begin{equation}
		\max \sum_{i=1}^{6}\sum_{j=1}^{7} S_{ij}(P_{ij} - B_i \times L_i)
		\label{eq:objective}
	\end{equation}
	
	约束条件：
	
	\begin{equation}
		R_{ij} \leq S_{imax}, \quad i=1,2,\dots,6; j=1,2,\dots,7
		\label{eq:constraint1}
	\end{equation}
	
	\begin{equation}
		110\% \times B_i \leq P_{ij} \leq 150\% \times B_i
		\label{eq:constraint2}
	\end{equation}
	
	求解得到未来一周最大收益为5105.60元。
	
	\newpage
	
	% ==================== 八、问题三模型的建立与求解 ====================
	\section{问题三模型的建立与求解}
	
	目标函数：
	
	\begin{equation}
		\max \sum_{i=1}^{n} y_i \left[ R_i(P_i - B_i) - R_i B_i L_i \right]
		\label{eq:obj3}
	\end{equation}
	
	约束条件：
	
	\begin{equation}
		27 \leq \sum_{i=1}^{n} y_i \leq 33
		\label{eq:const3_1}
	\end{equation}
	
	\begin{equation}
		R_i \geq 2.5, \quad i=1,2,\dots,n
		\label{eq:const3_2}
	\end{equation}
	
	\begin{equation}
		110\% \times B_i \leq P_i \leq 150\% \times B_i
		\label{eq:const3_3}
	\end{equation}
	
	求解得到2023年7月1日最大收益为1282.26元。
	
	\newpage
	
	% ==================== 九、问题四模型的建立与求解 ====================
	\section{问题四模型的建立与求解}
	
	灰色关联分析：
	
	\begin{equation}
		Z_{ij} = \frac{x_{ij}}{\bar{x}_{ij}}
		\label{eq:normalize}
	\end{equation}
	
	\begin{equation}
		\gamma(x_0(k), x_i(k)) = \frac{a + \rho b}{|x_0(k) - x_i(k)| + \rho b}
		\label{eq:grey_corr}
	\end{equation}
	
	\begin{table}[H]
		\centering
		\caption{三个指标与六大蔬菜品类季度平均销售量的关联度}
		\begin{tabular}{c|c|c|c}
			\hline
			\textbf{} & \textbf{指标1} & \textbf{指标2} & \textbf{指标3} \\
			\hline
			花菜类 & 0.9017 & 0.6937 & 0.6642 \\
			食用菌 & 0.8590 & 0.7287 & 0.7573 \\
			花叶类 & 0.7377 & 0.5961 & 0.5665 \\
			辣椒类 & 0.7219 & 0.5782 & 0.6297 \\
			\hline
		\end{tabular}
	\end{table}
	
	各指标与六大蔬菜品类季度平均销售量的关联度均在0.5以上，表明具有很强的相关性。
	
	\newpage
	
	% ==================== 十、模型的评价与推广 ====================
	\section{模型的评价与推广}
	
	\subsection{模型的优点}
	
	\begin{enumerate}
		\item K-means++算法是K-means算法的升级，在初始聚类中心上做出规定，选取了更合理的初始聚类中心。
		\item 线性回归模型可直观显示各蔬菜品类销售量的变化对成本加成定价的影响。
		\item 模拟退火算法可根据一定概率放弃局部最优解，找到全局最优解。
		\item 灰色关联分析简单易懂，可直观反映因素之间的关联程度。
	\end{enumerate}
	
	\subsection{模型的缺点}
	
	\begin{enumerate}
		\item K-means++算法初始中心选择不当，仍然可能陷入局部最优解。
		\item 灰色关联分析对数据要求过高，数据若存在较大异常值，可能会对结果产生较大影响。
	\end{enumerate}
	
	\subsection{模型的推广}
	
	时间序列可提取数据的长期、季节、循环和不规则趋势，可用于冷饮销售预测、服饰销售预测、天气温度预测等场景。
	
	\newpage
	
	% ==================== 十一、参考文献 ====================
	\section{参考文献}
	
	\begin{enumerate}
		\item 曾敏敏. 基于时间情境的生鲜社区超市的动态定价策略研究[D]. 2021.
		\item 魏泽娥, 陈刚, 丁胜春, 耿军霞. 大规模定制产品的成本加成定价方法研究[J]. 黑龙江对外经贸, 2007, (02): 84-85.
		\item 刘保政, 刘德宝, 高立群. 供不应求季节性商品的价格控制和生产销售决策模型[J]. 东北大学学报, 2005(11): 23-26.
		\item 王艳. 中小型连锁超市定价的影响因素及其定价策略探析[J]. 兰州工业学院学报, 2014, 21(03): 99-102.
	\end{enumerate}
	
	\newpage
	
\newpage

	% ==================== 附录 ====================
	\section{附录}
	
	\subsection{支撑材料文件列表}
	
	\begin{table}[H]
		\centering
		\begin{tabular}{c|c}
			\hline
			\textbf{文件分类} & \textbf{文件内容} \\
			\hline
			问题一 & 描述统计、Spearman相关系数、K-means++聚类 \\
			问题二 & 线性回归、时间序列预测、优化模型 \\
			问题三 & 模拟退火算法 \\
			问题四 & 灰色关联分析 \\
			\hline
		\end{tabular}
	\end{table}
	
	\subsection{单品聚类结果（完整）}
	
	\begin{longtable}{c|c|c}
		\caption{单品蔬菜聚类结果（完整）} \\
		\hline
		\textbf{序号} & \textbf{名称} & \textbf{类别} \\
		\hline
		\endfirsthead
		\hline
		\textbf{序号} & \textbf{名称} & \textbf{类别} \\
		\hline
		\endhead
		\hline
		\multicolumn{3}{r}{（接上页）} \\
		\endfoot
		\hline
		\endlastfoot
		
		% 热销类
		1 & 金针菇 & 热销 \\
		2 & 净藕 & 热销 \\
		3 & 芜湖青椒 & 热销 \\
		4 & 西兰花 & 热销 \\
		5 & 云南生菜 & 热销 \\
		\hline
		% 畅销类
		6 & 菠菜 & 畅销 \\
		7 & 大白菜 & 畅销 \\
		8 & 螺丝椒 & 畅销 \\
		9 & 奶白菜 & 畅销 \\
		10 & 上海青 & 畅销 \\
		11 & 西峡香菇 & 畅销 \\
		12 & 小米椒 & 畅销 \\
		13 & 云南油麦菜 & 畅销 \\
		14 & 紫茄子 & 畅销 \\
		\hline
		% 平销类
		15 & 白玉菇 & 平销 \\
		16 & 保康高山大白菜 & 平销 \\
		17 & 菜心 & 平销 \\
		18 & 海鲜菇 & 平销 \\
		19 & 红椒 & 平销 \\
		20 & 红薯尖 & 平销 \\
		21 & 洪湖莲藕（粉藕） & 平销 \\
		22 & 黄白菜 & 平销 \\
		23 & 黄心菜 & 平销 \\
		24 & 牛首油菜 & 平销 \\
		25 & 泡泡椒（精品） & 平销 \\
		26 & 青梗散花 & 平销 \\
		27 & 青茄子 & 平销 \\
		28 & 青线椒 & 平销 \\
		29 & 双孢菇 & 平销 \\
		30 & 甜白菜 & 平销 \\
		31 & 茼蒿 & 平销 \\
		32 & 娃娃菜 & 平销 \\
		33 & 苋菜 & 平销 \\
		34 & 小青菜 & 平销 \\
		35 & 小皱皮 & 平销 \\
		36 & 杏鲍菇 & 平销 \\
		37 & 枝江红菜苔 & 平销 \\
		38 & 枝江青梗散花 & 平销 \\
		39 & 竹叶菜 & 平销 \\
		\hline
		% 滞销类
		40 & 艾蒿 & 滞销 \\
		41 & 白菜苔 & 滞销 \\
		42 & 白蒿 & 滞销 \\
		43 & 薄荷叶 & 滞销 \\
		44 & 本地黄心油菜 & 滞销 \\
		45 & 本地上海青 & 滞销 \\
		46 & 本地小毛白菜 & 滞销 \\
		47 & 荸荠 & 滞销 \\
		48 & 冰草 & 滞销 \\
		49 & 蔡甸藜蒿 & 滞销 \\
		50 & 茶树菇（袋） & 滞销 \\
		51 & 赤松茸 & 滞销 \\
		52 & 虫草花 & 滞销 \\
		53 & 春菜 & 滞销 \\
		54 & 大白菜秧 & 滞销 \\
		55 & 大芥兰 & 滞销 \\
		56 & 大龙茄子 & 滞销 \\
		57 & 灯笼椒 & 滞销 \\
		58 & 东门口小白菜 & 滞销 \\
		59 & 甘蓝叶 & 滞销 \\
		60 & 高瓜 & 滞销 \\
		61 & 黑牛肝菌 & 滞销 \\
		62 & 黑皮鸡枞菌 & 滞销 \\
		63 & 黑油菜 & 滞销 \\
		64 & 红灯笼椒 & 滞销 \\
		65 & 红枞椒 & 滞销 \\
		66 & 红尖椒 & 滞销 \\
		67 & 红莲藕带 & 滞销 \\
		68 & 红珊瑚（粗叶） & 滞销 \\
		69 & 红线椒 & 滞销 \\
		70 & 红橡叶 & 滞销 \\
		71 & 洪湖莲藕（脆藕） & 滞销 \\
		72 & 洪湖藕带 & 滞销 \\
		73 & 洪山菜苔 & 滞销 \\
		74 & 猴头菇 & 滞销 \\
		75 & 花菇（一人份） & 滞销 \\
		76 & 花茄子 & 滞销 \\
		77 & 槐花 & 滞销 \\
		78 & 黄花菜 & 滞销 \\
		79 & 活体银耳 & 滞销 \\
		80 & 鸡枞菌 & 滞销 \\
		81 & 姬菇 & 滞销 \\
		82 & 荠菜 & 滞销 \\
		83 & 芥菜 & 滞销 \\
		84 & 芥兰 & 滞销 \\
		85 & 菊花油菜 & 滞销 \\
		86 & 快菜 & 滞销 \\
		87 & 辣妹子 & 滞销 \\
		88 & 莲蓬（个） & 滞销 \\
		89 & 菱角 & 滞销 \\
		90 & 龙牙菜 & 滞销 \\
		91 & 鹿茸菇（盒） & 滞销 \\
		92 & 萝卜叶 & 滞销 \\
		93 & 绿牛油 & 滞销 \\
		94 & 马齿苋 & 滞销 \\
		95 & 马兰头 & 滞销 \\
		96 & 面条菜 & 滞销 \\
		97 & 木耳菜 & 滞销 \\
		98 & 奶白菜苗 & 滞销 \\
		99 & 南瓜尖 & 滞销 \\
		100 & 牛排菇 & 滞销 \\
		101 & 牛首生菜 & 滞销 \\
		102 & 藕尖 & 滞销 \\
		103 & 平菇 & 滞销 \\
		104 & 蒲公英 & 滞销 \\
		105 & 七彩椒 & 滞销 \\
		106 & 青菜苔 & 滞销 \\
		107 & 青枞椒 & 滞销 \\
		108 & 青尖椒 & 滞销 \\
		109 & 双沟白菜 & 滞销 \\
		110 & 水果辣椒 & 滞销 \\
		111 & 丝瓜尖 & 滞销 \\
		112 & 四川红香椿 & 滞销 \\
		113 & 随州泡泡青 & 滞销 \\
		114 & 田七 & 滞销 \\
		115 & 豌豆尖 & 滞销 \\
		116 & 西峡花菇 & 滞销 \\
		117 & 鲜木耳 & 滞销 \\
		118 & 鲜藕带（袋） & 滞销 \\
		119 & 鲜粽叶 & 滞销 \\
		120 & 襄甜红菜苔（袋） & 滞销 \\
		121 & 小白菜 & 滞销 \\
		122 & 小鲜味菇 & 滞销 \\
		123 & 小炒菇 & 滞销 \\
		124 & 绣球菌 & 滞销 \\
		125 & 野藕 & 滞销 \\
		126 & 野生粉藕 & 滞销 \\
		127 & 银耳（朵） & 滞销 \\
		128 & 油菜苔 & 滞销 \\
		129 & 余干椒 & 滞销 \\
		130 & 鱼腥草 & 滞销 \\
		131 & 圆茄子 & 滞销 \\
		132 & 长线茄 & 滞销 \\
		133 & 芝麻苋菜 & 滞销 \\
		134 & 猪肚菇（盒） & 滞销 \\
		135 & 紫白菜 & 滞销 \\
		136 & 紫贝菜 & 滞销 \\
		137 & 紫尖椒 & 滞销 \\
		138 & 紫螺丝椒 & 滞销 \\
		139 & 紫苏 & 滞销 \\
		140 & 紫圆茄 & 滞销 \\
		
	\end{longtable}
	
	\newpage
	
	\subsection{核心代码}
	
	\begin{lstlisting}[language=Matlab, basicstyle=\small\ttfamily]
		% K-means++ clustering
		X = xlsread('data.xlsx', 'sheet1', 'B2:D41');
		K = 4;
		max_iters = 20;
		centroids = init_centroids(X, K);
		
		for i = 1:max_iters
		labels = assign_labels(X, centroids);
		centroids = update_centroids(X, labels, K);
		end
		
		function labels = assign_labels(X, centroids)
		[~, labels] = min(pdist2(X, centroids, 'squaredeuclidean'), [], 2);
		end
	\end{lstlisting}
	
	\begin{lstlisting}[language=Matlab, basicstyle=\small\ttfamily]
		% Simulated Annealing
		maxgen = 100;
		Lk = 100;
		T = 100;
		alfa = 0.95;
		
		for iter = 1:maxgen
		for i = 1:Lk
		x_new = x0 + 2 * randn(1, narvs);
		y_new = Obj_fun2(x_new);
		if y_new > y0 || exp(-(y0 - y_new) / T) > rand(1)
		x0 = x_new;
		y0 = y_new;
		end
		end
		T = alfa * T;
		end
	\end{lstlisting}
	
	\subsection{单品蔬菜每日批发单价与损耗率表}
	
	\begin{longtable}{c|c|c|c|c|c}
		\caption{单品蔬菜每日批发单价与损耗率表} \\
		\hline
		\textbf{日期} & \textbf{白玉菇（袋）} & \textbf{菠菜} & \textbf{菠菜（份）} & \textbf{菜心} & \textbf{虫草花（份）} \\
		\hline
		\endfirsthead
		\hline
		\textbf{日期} & \textbf{白玉菇（袋）} & \textbf{菠菜} & \textbf{菠菜（份）} & \textbf{菜心} & \textbf{虫草花（份）} \\
		\hline
		\endhead
		\hline
		\multicolumn{6}{r}{（接上页）} \\
		\endfoot
		\hline
		\endlastfoot
		
		2023-06-24 & 0 & 9.63 & 3.94 & 4.62 & 2.66 \\
		2023-06-25 & 0 & 9.63 & 0 & 4.62 & 2.66 \\
		2023-06-26 & 3.58 & 9.66 & 0 & 0 & 2.66 \\
		2023-06-27 & 0 & 9.67 & 4.13 & 4.61 & 2.67 \\
		2023-06-28 & 3.57 & 9.66 & 4.19 & 4.62 & 2.72 \\
		2023-06-29 & 3.57 & 0 & 4.20 & 2.63 & 2.71 \\
		2023-06-30 & 3.29 & 9.67 & 4.07 & 0 & 2.61 \\
		
	\end{longtable}
	
	\begin{longtable}{c|c|c|c|c|c}
		\caption{单品蔬菜每日批发单价与损耗率表（续）} \\
		\hline
		\textbf{日期} & \textbf{高瓜（1）} & \textbf{高瓜（2）} & \textbf{海鲜菇（包）} & \textbf{紫茄子（2）} & \textbf{红椒（2）} \\
		\hline
		\endfirsthead
		\hline
		\textbf{日期} & \textbf{高瓜（1）} & \textbf{高瓜（2）} & \textbf{海鲜菇（包）} & \textbf{紫茄子（2）} & \textbf{红椒（2）} \\
		\hline
		\endhead
		\hline
		\multicolumn{6}{r}{（接上页）} \\
		\endfoot
		\hline
		\endlastfoot
		
		2023-06-24 & 11.14 & 0 & 1.95 & 3.82 & 13.46 \\
		2023-06-25 & 11.31 & 3.23 & 1.96 & 3.82 & 12.80 \\
		2023-06-26 & 11.65 & 13.35 & 1.96 & 3.82 & 12.78 \\
		2023-06-27 & 11.64 & 13.39 & 1.95 & 3.82 & 12.91 \\
		2023-06-28 & 11.67 & 0 & 1.96 & 3.82 & 12.81 \\
		2023-06-29 & 11.71 & 13.67 & 1.95 & 3.82 & 12.80 \\
		2023-06-30 & 11.67 & 13.69 & 1.95 & 3.43 & 12.72 \\
		
	\end{longtable}
	
	\subsection{单品蔬菜补货量与定价策略}
	
	\begin{longtable}{c|c|c|c|c}
		\caption{单品蔬菜补货量与定价策略} \\
		\hline
		\textbf{序号} & \textbf{单品名称} & \textbf{单位成本} & \textbf{日补货总量（kg）} & \textbf{定价} \\
		\hline
		\endfirsthead
		\hline
		\textbf{序号} & \textbf{单品名称} & \textbf{单位成本} & \textbf{日补货总量（kg）} & \textbf{定价} \\
		\hline
		\endhead
		\hline
		\multicolumn{5}{r}{（接上页）} \\
		\endfoot
		\hline
		\endlastfoot
		
		1 & 本地小毛白菜 & 3.14 & 4.65 & 6.66 \\
		2 & 云南生菜 & 5.86 & 7.26 & 10.44 \\
		3 & 茼蒿 & 8.23 & 9.01 & 19.92 \\
		4 & 牛首油菜 & 2.84 & 18.22 & 5.50 \\
		5 & 紫茄子（2） & 4.56 & 13.89 & 6.67 \\
		6 & 西峡香菇（1） & 13.79 & 17.72 & 20.33 \\
		7 & 西兰花 & 10.04 & 18.04 & 16.07 \\
		8 & 羊角白 & 4.96 & 24.97 & 11.32 \\
		9 & 枝江黄花 & 5.64 & 2.50 & 9.01 \\
		10 & 白玉菇（袋） & 4.38 & 12.50 & 9.87 \\
		11 & 芜湖青椒（1） & 3.89 & 26.92 & 6.16 \\
		12 & 杏鲍菇（2） & 9.13 & 13.15 & 21.79 \\
		13 & 奶白菜（份） & 2.36 & 22.13 & 4.04 \\
		14 & 双孢菇（盒） & 3.48 & 10.53 & 5.32 \\
		15 & 青红杭椒组合装（份） & 3.53 & 5.62 & 10.47 \\
		16 & 红椒（2） & 14.12 & 2.76 & 28.54 \\
		17 & 蟹味菇与白玉菇双拼（盒） & 3.38 & 6.27 & 7.21 \\
		18 & 洪湖莲藕（粉藕） & 8.09 & 13.85 & 40.29 \\
		19 & 小青菜（1） & 3.07 & 5.32 & 6.96 \\
		20 & 云南生菜（份） & 2.24 & 13.00 & 6.54 \\
		21 & 云南油麦菜（份） & 2.80 & 9.00 & 5.44 \\
		22 & 菠菜（份） & 3.37 & 15.00 & 7.99 \\
		23 & 鱼腥草（份） & 1.91 & 2.50 & 3.12 \\
		24 & 小米椒（份） & 2.80 & 19.00 & 6.12 \\
		25 & 冰草（盒） & 5.56 & 13.00 & 12.14 \\
		26 & 金针菇（盒） & 1.69 & 20.00 & 2.96 \\
		27 & 黄白菜（2） & 5.07 & 7.80 & 9.95 \\
		28 & 圆茄子（2） & 4.59 & 4.02 & 6.38 \\
		29 & 奶白菜 & 2.79 & 6.47 & 3.80 \\
		
	\end{longtable}
	

	
\end{document}